\chapter{How M.~d'Eymeris's Diamond Passed into the Hands of M.~d'Artagnan}
	
	\chaptermark{How M.~d'Eymeris's Diamond Came To M.~d'Artagnan}

\lettrine{W}{hilst} this violent, noisy, and bloody scene was passing on the Grève, several men, barricaded behind the gate of communication with the garden, replaced their swords in their sheaths, assisted one among them to mount a ready saddled horse which was waiting in the garden, and like a flock of startled birds, fled in all directions, some climbing the walls, others rushing out at the gates with all the fury of a panic. He who mounted the horse, and gave him the spur so sharply that the animal was near leaping the wall, this cavalier, we say, crossed the Place Baudoyer, passed like lightening before the crowd in the streets, riding against, running over and knocking down all that came in his way, and, ten minutes after, arrived at the gates of the superintendent, more out of breath than his horse. The Abbé Fouquet, at the clatter of hoofs on the pavement, appeared at a window of the court, and before even the cavalier had set foot to the ground, <Well! Danicamp?> cried he, leaning half out of the window.

<Well, it is all over,> replied the cavalier.

<All over!> cried the Abbé. <Then they are saved?>

<No, monsieur,> replied the cavalier, <they are hung.>

<Hung!> repeated the Abbé, turning pale. A lateral door suddenly opened, and Fouquet appeared in the chamber, pale, distracted, with lips half opened, breathing a cry of grief and anger. He stopped upon the threshold to listen to what was addressed from the court to the window.

<Miserable wretches!> said the Abbé, <you did not fight, then?>

<Like lions.>

<Say like cowards.>

<Monsieur!>

<A hundred men accustomed to war, sword in hand, are worth ten thousand archers in a surprise. Where is Menneville, that boaster, that braggart, who was to come back either dead or a conqueror?>

<Well, monsieur, he kept his word. He is dead!>

<Dead! Who killed him?>

<A demon disguised as a man, a giant armed with ten flaming swords—a madman, who at one blow extinguished the fire, put down the riot, and caused a hundred musketeers to rise up out of the pavement of the Grève.>

Fouquet raised his brow, streaming with sweat, murmuring, <Oh! Lyodot and D'Eymeris! dead! dead! dead! and I dishonoured.>

The Abbé turned round, and perceiving his brother, despairing and livid, <Come, come,> said he, <it is a blow of fate, monsieur; we must not lament thus. Our attempt has failed because God\longdash>

<Be silent, Abbé! be silent!> cried Fouquet; <your excuses are blasphemies. Order that man up here, and let him relate the details of this terrible event.>

<But, brother\longdash>

<Obey, monsieur!>

The Abbé made a sign, and in half a minute the man's step was heard upon the stairs. At the same time Gourville appeared behind Fouquet, like the guardian angel of the superintendent, pressing one finger on his lips to enjoin observation even amidst the bursts of his grief. The minister resumed all the serenity that human strength left at the disposal of a heart half broken with sorrow. Danicamp appeared. <Make your report,> said Gourville.

<Monsieur,> replied the messenger, <we received orders to carry off the prisoners, and to cry <Vive Colbert!> whilst carrying them off.>

<To burn them alive, was it not, Abbé?> interrupted Gourville.

<Yes, yes, the order was given to Menneville. Menneville knew what was to be done, and Menneville is dead.>

This news appeared rather to reassure Gourville than to sadden him.

<Yes, certainly to burn them alive,> said the Abbé, eagerly.

<Granted, monsieur, granted,> said the man, looking into the eyes and the faces of the two interlocutors, to ascertain what there was profitable or disadvantageous to himself in telling the truth.

<Now, proceed,> said Gourville.

<The prisoners,> cried Danicamp, <were brought to the Grève, and the people, in a fury, insisted upon their being burnt instead of being hung.>

<And the people were right,> said the Abbé. <Go on.>

<But,> resumed the man, <at the moment the archers were broken, at the moment the fire was set to one of the houses of the Place destined to serve as a funeral-pile for the guilty, this fury, this demon, this giant of whom I told you, and who, we had been informed, was the proprietor of the house in question, aided by a young man who accompanied him, threw out of the window those who kept the fire, called to his assistance the musketeers who were in the crowd, leaped himself from the window of the first story into the Place, and plied his sword so desperately that the victory was restored to the archers, the prisoners were retaken, and Menneville killed. When once recaptured, the condemned were executed in three minutes.> Fouquet, in spite of his self-command, could not prevent a deep groan escaping him.

<And this man, the proprietor of the house, what is his name?> said the Abbé.

<I cannot tell you, not having even been able to get sight of him; my post had been appointed in the garden, and I remained at my post: only the affair was related to me as I repeat it. I was ordered, when once the affair was at an end, to come at best speed and announce to you the manner in which it finished. According to this order, I set out, full gallop, and here I am.>

<Very well, monsieur, we have nothing else to ask of you,> said the Abbé, more and more dejected, in proportion as the moment approached for finding himself alone with his brother.

<Have you been paid?> asked Gourville.

<Partly, monsieur,> replied Danicamp.

<Here are twenty pistols. Begone, monsieur, and never forget to defend, as this time has been done, the true interests of the king.>

<Yes, monsieur,> said the man, bowing and pocketing the money. After which he went out. Scarcely had the door closed after him when Fouquet, who had remained motionless, advanced with a rapid step and stood between the Abbé and Gourville. Both of them at the same time opened their mouths to speak to him. <No excuses,> said he, <no recriminations against anybody. If I had not been a false friend I should not have confided to any one the care of delivering Lyodot and D'Eymeris. I alone am guilty; to me alone are reproaches and remorse due. Leave me, Abbé.>

<And yet, monsieur, you will not prevent me,> replied the latter, <from endeavouring to find out the miserable fellow who has intervened to the advantage of M.~Colbert in this so well-arranged affair; for, if it is good policy to love our friends dearly, I do not believe that is bad which consists in obstinately pursuing our enemies.>

<A truce to policy, Abbé; begone, I beg of you, and do not let me hear any more of you till I send for you; what we most need is circumspection and silence. You have a terrible example before you, gentlemen: no reprisals, I forbid them.>

<There are no orders,> grumbled the Abbé, <which will prevent me from avenging a family affront upon the guilty person.>

<And I,> cried Fouquet, in that imperative tone to which one feels there is nothing to reply, <if you entertain one thought, one single thought, which is not the absolute expression of my will, I will have you cast into the Bastille two hours after that thought has manifested itself. Regulate your conduct accordingly, Abbé.>

The Abbé coloured and bowed. Fouquet made a sign to Gourville to follow him, and was already directing his steps towards his cabinet, when the usher announced with a loud voice: <Monsieur le Chevalier d'Artagnan.>

<Who is he?> said Fouquet, negligently, to Gourville.

<An ex-lieutenant of his majesty's musketeers,> replied Gourville, in the same tone. Fouquet did not even take the trouble to reflect, and resumed his walk. <I beg your pardon, monseigneur!> said Gourville, <but I have remembered; this brave man has quitted the king's service, and probably comes to receive an instalment of some pension or other.>

<Devil take him!> said Fouquet, <why does he choose his opportunity so ill?>

<Permit me then, monseigneur, to announce your refusal to him; for he is one of my acquaintance, and is a man whom, in our present circumstances, it would be better to have as a friend than an enemy.>

<Answer him as you please,> said Fouquet.

<Eh! good Lord!> said the Abbé, still full of malice, like an egotistical man; <tell him there is no money, particularly for musketeers.>

But scarcely had the Abbé uttered this imprudent speech, when the partly open door was thrown back, and D'Artagnan appeared.

<Eh! Monsieur Fouquet,> said he, <I was well aware there was no money for musketeers here. Therefore I did not come to obtain any, but to have it refused. That being done, receive my thanks. I give you good-day, and will go and seek it at M.~Colbert's.> And he went out, making an easy bow.

<Gourville,> said Fouquet, <run after that man and bring him back.> Gourville obeyed, and overtook D'Artagnan on the stairs.

D'Artagnan, hearing steps behind him, turned round and perceived Gourville. <\swear{Mordioux!} my dear monsieur,> said he, <there are sad lessons which you gentlemen of finance teach us; I come to M.~Fouquet to receive a sum accorded by his majesty, and I am received like a mendicant who comes to ask charity, or a thief who comes to steal a piece of plate.>

<But you pronounced the name of M.~Colbert, my dear M.~d'Artagnan; you said you were going to M.~Colbert's?>

<I certainly am going there, were it only to ask satisfaction of the people who try to burn houses, crying <Vive Colbert!>>

Gourville pricked up his ears. <Oh, oh!> said he, <you allude to what has just happened at the Grève?>

<Yes, certainly.>

<And in what did that which has taken place concern you?>

<What! do you ask me whether it concerns me or does not concern me, if M.~ Colbert pleases to make a funeral-pile of my house?>

<So, ho, your house—was it your house they wanted to burn?>

<\swear{Pardieu!} was it!>

<Is the cabaret of the Image-de-Notre-Dame yours, then?>

<It has been this week.>

<Well, then, are you the brave captain, are you the valiant blade who dispersed those who wished to burn the condemned?>

<My dear Monsieur Gourville, put yourself in my place. I was an agent of the public force and a landlord, too. As a captain, it is my duty to have the orders of the king accomplished. As a proprietor, it is to my interest my house should not be burnt. I have at the same time attended to the laws of interest and duty in replacing Messieurs Lyodot and D'Eymeris in the hands of the archers.>

<Then it was you who threw the man out of the window?>

<It was I, myself,> replied D'Artagnan, modestly.

<And you who killed Menneville?>

<I had that misfortune,> said D'Artagnan, bowing like a man who is being congratulated.

<It was you, then, in short, who caused the two condemned persons to be hung?>

<Instead of being burnt, yes, monsieur, and I am proud of it. I saved the poor devils from horrible tortures. Understand, my dear Monsieur de Gourville, that they wanted to burn them alive. It exceeds imagination!>

<Go, my dear Monsieur d'Artagnan, go,> said Gourville, anxious to spare Fouquet the sight of the man who had just caused him such profound grief.

<No,> said Fouquet, who had heard all from the door of the ante-chamber; <not so; on the contrary, Monsieur d'Artagnan, come in.>

D'Artagnan wiped from the hilt of his sword a last bloody trace, which had escaped his notice, and returned. He then found himself face to face with these three men, whose countenances wore very different expressions. With the Abbé it was anger, with Gourville stupor, with Fouquet it was dejection.

<I beg your pardon, monsieur le ministre,> said D'Artagnan, <but my time is short; I have to go to the office of the intendant, to have an explanation with Monsieur Colbert, and to receive my quarter's pension.>

<But, monsieur,> said Fouquet, <there is money here.> D'Artagnan looked at the superintendent with astonishment. <You have been answered inconsiderately, monsieur, I know, because I heard it,> said the minister; <a man of your merit ought to be known by everybody.> D'Artagnan bowed. <Have you an order?> added Fouquet.

<Yes, monsieur.>

<Give it me, I will pay you myself; come with me.> He made a sign to Gourville and the Abbé, who remained in the chamber where they were. He led D'Artagnan into his cabinet. As soon as the door was shut,—<how much is due to you, monsieur?>

<Why, something like five thousand livres, monseigneur.>

<For arrears of pay?>

<For a quarter's pay.>

<A quarter consisting of five thousand livres!> said Fouquet, fixing upon the musketeer a searching look. <Does the king, then, give you twenty thousand livres a year?>

<Yes, monseigneur, twenty thousand livres a year. Do you think it is too much?>

<I\@?> cried Fouquet, and he smiled bitterly. <If I had any knowledge of mankind, if I were—instead of being a frivolous, inconsequent, and vain spirit—of a prudent and reflective spirit; if, in a word, I had, as certain persons have known how, regulated my life, you would not receive twenty thousand livres a year, but a hundred thousand, and you would belong not to the king but to me.>

D'Artagnan coloured slightly. There is sometimes in the manner in which a eulogium is given, in the voice, in the affectionate tone, a poison so sweet, that the strongest mind is intoxicated by it. The superintendent terminated his speech by opening a drawer, and taking from it four rouleaux, which he placed before D'Artagnan. The Gascon opened one. <Gold!> said he.

<It will be less burdensome, monsieur.>

<But, then, monsieur, these make twenty thousand livres.>

<No doubt they do.>

<But only five are due to me.>

<I wish to spare you the trouble of coming four times to my office.>

<You overwhelm me, monsieur.>

<I do only what I ought to do, monsieur le chevalier; and I hope you will not bear me any malice on account of the rude reception my brother gave you. He is of a sour, capricious disposition.>

<Monsieur,> said D'Artagnan, <believe me, nothing would grieve me more than an excuse from you.>

<Therefore I will make no more, and will content myself with asking you a favour.>

<Oh, monsieur.>

Fouquet drew from his finger a ring worth about three thousand pistoles. <Monsieur,> said he, <this stone was given me by a friend of my childhood, by a man to whom you have rendered a great service.>

<A service—I\@?> said the musketeer; <I have rendered a service to one of your friends?>

<You cannot have forgotten it, monsieur, for it dates this very day.>

<And that friend's name was\longdash>

<M.~d'Eymeris.>

<One of the condemned?>

<Yes, one of the victims. Well! Monsieur d'Artagnan, in return for the service you have rendered him, I beg you to accept this diamond. Do so for my sake.>

<Monsieur! you\longdash>

<Accept it, I say. To-day is with me a day of mourning; hereafter you will, perhaps, learn why; to-day I have lost one friend; well, I will try to get another.>

<But, Monsieur Fouquet\longdash>

<Adieu! Monsieur d'Artagnan, adieu!> cried Fouquet, with much emotion; <or rather, au revoir.> And the minister quitted the cabinet, leaving in the hands of the musketeer the ring and the twenty thousand livres.

<Oh!> said D'Artagnan, after a moment's dark reflection. <How on earth am I to understand what this means? \swear{Mordioux!} I can understand this much, only: he is a gallant man! I will go and explain matters to M.~Colbert.> And he went out. 